While the integration of LLMs and agents has shown strong potential for automating kernel generation, the field remains at an early stage of development. Bridging the gap between promising prototypes and production-grade systems requires addressing a set of interrelated challenges. This section examines these challenges and highlights emerging research directions spanning data, agents, infrastructure, evaluation, and human–AI collaboration, which are likely to shape the next generation of AI-driven kernel generation and optimization systems.


% \label{sec:data_scarcity}

\paragraph{Data Scarcity and Synthetic Scaling.} Progress toward production-grade performance remains fundamentally constrained by data scarcity. High-performance kernels exhibit a pronounced long-tail distribution and are sparsely represented in existing code corpora, where most available datasets still lack deep, hardware-aware domain knowledge, and existing corpora predominantly capture only final optimized kernels but omit optimization trajectories. Promising directions to access these limitations include systematic kernel dataset construction, large-scale synthetic data generation, and the collection of execution-driven optimization processes. Such data can support a wide range of learning paradigms, including pretraining, supervised fine-tuning, and reinforcement learning, and may be crucial for enabling meaningful scaling behavior in kernel generation systems.

% \subsection{}

\paragraph{Agentic Reasoning and Engineering Standards}. Current agent-based kernel optimization relies on predefined, workflow-driven paradigms, often failing long-horizon tasks due to redundant exploration and context exhaustion. To transcend these limitations, we propose three critical advancements: (1) enhancing autonomy by shifting from handcrafted workflows to self-directed planning and dynamic memory; (2) enabling principled reasoning by integrating dispersed heuristics across documentation and experts into structured knowledge bases; and (3) ensuring reliability through rigorous engineering standards, including formal verification and strict specifications. Collectively, addressing these challenges is critical for transforming agentic kernel optimization from exploratory automation into a robust, engineering-grade capability.


% \subsection{}

\paragraph{Scalable Infrastructure for Synthesis and Training.} Scalable infrastructure remains a bottleneck due to the severe latency mismatch between rapid model inference and costly kernel compilation. This disparity hinders the high-throughput feedback loops essential for reinforcement learning and synthetic data generation. Addressing this challenge calls for infrastructure that cleanly decouples model reasoning from environment execution via standardized, distributed ``gym-like'' environments, while supporting distributed and asynchronous execution at scale. Ultimately, advances in scalable infrastructure are critical for transforming kernel synthesis and data sampling from low-throughput experimentation into a systematic, data-driven learning process.


% \subsection{}

% A key open challenge in AI-driven kernel generation is the lack of robust and comprehensive evaluation. Many existing benchmarks assess performance under narrowly defined settings, limiting confidence in reported gains and obscuring generalization behavior. 

% At the operator level, evaluations are often restricted to fixed input shapes, limited precision regimes, and a small set of forward-pass primitives. Such protocols fail to reflect real-world workloads, where kernels must generalize across varying shapes, data types, and execution contexts. Moreover, backward operators, fused kernels, and composite workloads remain underrepresented, raising concerns about the applicability of current methods beyond canonical benchmarks. Evaluation limitations are further compounded by ecosystem concentration. Most approaches are developed and validated within the NVIDIA ecosystem, predominantly targeting CUDA-based kernels. This focus leaves open the question of whether current methodological paradigms such as data synthesis, agentic optimization, and reinforcement learning—can be instantiated across diverse hardware platforms, programming models, and toolchains. Addressing these gaps requires evaluation protocols that jointly assess robustness and generalization across shapes, operators, and ecosystems, providing a more reliable foundation for measuring progress in kernel synthesis research.

\paragraph{Evaluation Robustness and Generalization.} A key open challenge in AI-driven kernel generation is the lack of robust and comprehensive evaluation. A critical deficit in AI-driven kernel generation is the lack of robust evaluation. Existing benchmarks are often confined to fixed input shapes and forward-pass primitives within the NVIDIA ecosystem, failing to reflect the diversity of real-world workloads.  Addressing these gaps requires evaluation protocols that jointly assess robustness and generalization across shapes, operators, and ecosystems, providing a more reliable foundation for measuring progress in kernel generation research.

% \subsection{}

% Beyond fully automated approaches, human–AI collaboration represents an important and complementary paradigm for kernel generation. An open research question is how to systematically combine agentic exploration with human expertise to expand the design space and improve controllability in performance-critical settings.

% A central challenge in this paradigm is explainability. To enable effective collaboration and oversight, agents must provide interpretable rationales for their optimization decisions such as tiling strategies and memory layouts, which allow humans to reason about, validate, and refine system behavior. Without transparent explanations, human involvement is reduced to ad-hoc intervention, limiting the effectiveness of collaborative workflows. Equally important is the design of interaction paradigms that support principled division of labor between humans and agents. Future systems should facilitate mixed-initiative collaboration, in which humans specify high-level intent, constraints, or optimization objectives, while agents perform implementation, exploration, and parameter tuning. Establishing such interaction models is essential for enabling flexible and scalable collaboration without sacrificing automation benefits.


\paragraph{Human-AI Collaboration for Kernel Generation.} Beyond fully automated approaches, human–AI collaboration represents an important and complementary paradigm for kernel generation. An open research question is how to systematically combine agentic exploration with human expertise to expand the design space and improve controllability in performance-critical settings. To operationalize this, we identify two critical requirements: (1) Explainability, where agents provide interpretable rationales for optimization decisions (e.g., tiling) to facilitate expert verification; and (2) Mixed-initiative interaction, a paradigm where humans specify high-level constraints while agents execute implementation and tuning. Establishing this principled division of labor is essential to balance controllability with the scalability of automation.